% ! TeX root = distributed-systems-final-report.tex
\section{Requirements Elicitation and Analysis}\label{requirements}

Generally, the functional requirements of the project are to build the said
simulation, together with an application to visualize it and some commands to
manage it.

\begin{itemize}
  \item 1 - The system shall comprise a roundabout, surrounded by a certain
  number of straight roads (e.g. 4).
  \item 1.1 - The roundabout has a single lane and vehicles inside move
  counter-clockwise.
  \item 1.2 - A road is made of two lanes, for opposite directions; vehicles
  ride on their respective right side.
  \item 2 - The system shall comprise a certain amount of cars (in a reasonable
  number, e.g. up to 30).
  \item 2.1 - A car's surface covers a rectangle, both visually and
  physically (e.g. for collisions).
  \item 2.2 - A car is able to move in space, at a speed which can range
  from a certain minimum to a certain maximum and be changed by changing its
  acceleration. The physical model thus comprehends simple uniform motion
  and uniformly accelerated motion, both linear (for a straight road) and
  circular (around the roundabout).
  \item 2.3 - Each car is assigned an entry road and an exit road: it
  will come from the former towards the central roundabout, ride some distance
  inside it and then exit from the latter. It can spawn on any point on the
  entry road and despawn once it's ridden for some distance on the exit road.
  The distance ridden inside the roundabout must be strictly greater than zero
  and no more than a full lap, meaning that having the same entry and exit road
  implicates doing exactly one lap.
  \item 2.4 - Each car is able to proceed autonomously.
  \item 2.5 - Each car is provided with a local vision, allowing it to perceive
  others nearby (in reality, it could be implemented through infra-red, for
  example).t
  \item 3 - The system shall comprise one controller.
  \item 3.1 - The controller can receive spatial data from vehicles.
  \item 3.2 - The controller can send commands to vehicles, to alter heir
  acceleration.
  \item 4 - The system shall show a simulation, where vehicles can move from
  each one's start to end point and communicate with the controller.
  \item 4.1 - In the simulation, vehicles can be in one of three states: normal
  (correctly communicating with the controller and following its orders);
  failsafe (correctly communicating with the controller, but not receiving, or
  receiving but not following, its orders); disconnected (completely unable to
  either send or receive data to or from the controller).
  \item 4.2 - The simulation constantly checks for collisions between cars, in
  order to report them.
  \item 5 - The system shall offer ways to both visualize and give some commands
  to the simulation.
  \item 5.1 - The system allows to visualize the simulation, in real time,
  through a web application.
  \item 5.2 - The web application allows to pause or resume the simulation.
  \item 5.3 - The web application allows to visualize additional data on the
  vehicles, such as current speed and acceleration.
  \item 5.4 - The system allows to send additional commands via http requests,
  also accessible through a \textit{bash} wrapper terminal script.
  \item 5.5 - Available commands include: pausing and resuming the simulation,
  changing the state of one or more vehicles.
\end{itemize}

Non-functional requirements are instead related to safety of vehicles on the
road:

\begin{itemize}
  \item 6.1\label{req:req_fault} - \textbf{Fault tolerance} of the entire
  system: the controller or the network may even experience malfunctions, but
  the safety of vehicles should be maintaned at all times, to prevent accidents.
  \item\label{req:req_openness} 6.2 - \textbf{Opennes} and
  \textbf{interoperability}: the system should be open and interoperable with
  different implementations, for example, from different car brands.
\end{itemize}

\todo{tdd/unit tests/other tests; model first etc.}


\subsection{Relevant Distributed System Features}

\subsubsection{Transparency}\label{req:distr_transparency}
Despite the general advantages of a higher level of abstraction for the user,
\textbf{transparency} is only required in a few forms. The main one is
\textbf{scalability transparency}, since the number of vehicles will vary over
time (also refer to \ref{req:distr_scalability}). \textbf{Replication, migration}
and \textbf{performance transparency}, instead, aren't strictly required, but
are always a good feature for the central controller. Lastly, \textbf{location
transparency} is necessary, since vehicles are moving all the time.

On the other hand, \textbf{concurrency transparency} isn't relevant, since
there's no shared resource to access, and \textbf{failure transparency} isn't
granted, since failures are actually taken in great consideration in the project
requirements themselves.


Several forms of \textbf{transparency} can be found in this project, although
most aren't strongly required, if not for the general advantages of having a
higher level of abstraction for the user. The system should just be resilient
enough to accept a variable number of users, i.e. vehicles, replacing old ones with
new ones as they go away; also, the location of nodes isn't important and, given
the nature of the context itself, vehicles move all the times (\textbf{location
transparency}) - although they of course have to share their position.

\textbf{concurrency transparency} isn't relevant, since there's no shared
resource to access (see \ref{req:distr_sharing}). So, \textbf{scalability
transparency} is relevant for users, but it has not been handled for the central
controller, nor have been other forms like \textbf{replication, migration,
performance}, although the broker intermediary would help to hide these details.
Finally, \textbf{failure transparency} isn't granted, since both the controller
and a user will know if the other isn't responding, in which case the car will
switch to \textit{failsafe mode}.

\subsubsection{Dependability}\label{req:distr_dependability}
\textbf{Fault tolerance} is the main focus of the project, since \textbf{safety}
should be granted at all times. Even with network issues, the system should
provide \textbf{reliability} globally.


\textbf{Fault tolerance} is the main focus of the project, since \textbf{safety}
should be granted at all times. Even if the controller may not respond, the
entire system provides \textbf{reliability}, mainly through two measures
implemented on every single vehicle: there's no alternative to prevent all
accidents.

First, after a network timeout, a \textit{failsafe mode} is activated
immediately, providing instructions in the place of the controller without
disrupting the traffic flow. Thanks to it, users' nodes are totally autonomous
and independent, so that the controller isn't necessary but only increases the
efficiency of the system.

Secondly, even while remote commands are being received correctly, they can be
\textit{overridden} in any moment if there's an emergency: for example, the
vehicle may start to get too close to the one to the front, or even a totally
unexpected new vehicle may appear ahead, one which was disconnected from the
network and the controller wasn't aware about - besides other possible ostacles
in the real world, such as pedestrian and fallen objects. This measure could
also prove useful if communication data got somehow altered, either by
interferences or by malicious attackers.

Both these autonomous operation modes rely on the car's \textit{local vision}.
The \textbf{controller} also has an additional way to handle unresponding nodes:
besides vehicles' data, it also receives what appears in the local vision of
each one. Thanks to this sort of crowdsensing, it's able to reconstruct a map of
all vehicles, so that it can make its calculations as optimal as possible, even
if it won't be able to send commands to these so-called \textit{ghosts}. It
should be noted, however, that the more vehicles are disconnected from the
network, the less reliable this data will be, having more ``invisible'' nodes
than those contributing to the crowdsensing. Moreover, data can only be provided
by vehicles connected to the network (whether they're responding to commands or
driving in failsafe mode), but not by those which are completely disconnected
from the network.

\subsubsection{Scalability}\label{req:distr_scalability}
As already mentioned, the number of vehicles will vary over time; however, it
should be noted that the roundabout's surroundings can only contain a limited
amount of cars, hence the actual number of involved nodes at the same time isn't
expected to grow a lot. For a different, larger scenario trying to coordinate
traffic across a wider area, this may become a more urgent issue.


In addition to what said in \ref{req:distr_transparency}, \textbf{scaling} can
easily be performed thanks to the publish-subscribe model: no additional
configuration would be required for the communication with vehicles, although
increasing the amount of controller nodes would require adding some complexity
to the system. Moreover, the MQTT protocol is specifically designed to work with
constrained devices, to grant netowrk reliability. On the other hand, in its
current implementation, the system isn't expected to expand much: the roundabout
surrounding can only contain a limited amount of cars, and the old ones are
constantly replaced, in the controller's memory, as they drive away.

\subsubsection{Security}
\textbf{Security} would be very important in this scenario, in
the real world, since giving false information would mislead the controller's
calculations on one side and result in dangerous instructions for the vehicles
on the other: thus, it would be particularly important to guarantee data
\textbf{integrity} and use \textbf{authentication} to prevent the injection of
false data.

On one hand, in fact, wrong commands could lead to accidents; on the other hand,
an overflooding of inexisting cars' positions could lead the controller to block
the entire system.


\textbf{Security} wasn't taken into consideration in the development of this
project (for lack of time), but it would be very important in this scenario, in
the real world, since giving false information would mislead the controller's
calculations on one side and result in dangerous instructions for the vehicles
on the other: thus, it would be particularly important to guarantee data
\textbf{integrity} and use \textbf{authentication} to prevent the injection of
false data.
TLS authentication, already supported by Mosquitto MQTT, could easily be added
in the future.

According to the current design, the system has already been thought
to prevent accidents due to wrong commands, thanks to the cars' \textit{local
vision} and the \textit{failsafe mode} that can take control in case of
emergency; however, an overflooding of inexisting cars' positions could lead the
controller to block the entire system.

\subsubsection{Resource sharing}\label{req:distr_sharing}
No resource has been planned to be shared.


There's no shared resource to access, but rather each car shares its position
asynchronously and then the controller will periodically evaluate available data
and issue commands back, one by one, through MQTT messages.

\subsubsection{Openness, interoperability, heterogeneity of components}\label{req:distr_openness}
As stated in constraint 6.2 (\ref{req:req_openness}), the system should be open to
potentially welcome external vehicles' implementations, for example even with
different languages and technologies.


Using standardized communication protocols (JSON messages over MQTT), the system
is open to potentially welcome external vehicles' implementations, for example
even with different languages and technologies: any vehicle with similar
characteristics may join the network and adhere to the controller's supervision.

\subsubsection{Evolvability, maintainability}\ref{req:distr_evolv}
The system probably doesn't require updates, unless there are changes to the
road or to traffic rules, but it may be extended in the future with additional
features or performance optimizations. The logs could be useful to collect and
analyze, for example, in case of an accident.

\subsubsection{Performance, concurrency, computation/communmication efficiency, bandwidth}
Since vehicles are constantly moving, possibly even at high speeds, it's
important to pose a limit to the latency of the communication of commands from
the controller to each one.

Regarding computational performance and network load, since the system isn't
expected to grow significantly (as mentioned in \ref{req:distr_scalability}),
there shouldn't be a need for particularly expensive servers and network
interfaces to sustain enough vehicles.


Since vehicles are constantly moving, possibly even at high speeds, the command
evaluation frequency is important, although, as well as the timeout to consider
data too old (and fallback to \textit{failsafe mode}).

An higher number of cycles per seconds (currently 10 for both cars and
controller) would improve estimations, in exchange for computational cost (cars,
having a simpler logic, could easily run at higher frequency, but they were only
limited in the simulation for efficiency purposes). The performance impact would
even worsen with more users, although a not
too expensive server should be able to run the controller efficiently.

A higher network timeout (currently 0.5 seconds), instead, would result in more
wrong estimations from the controller and, thus, more activations of
\textit{failsafe mode} to prevent crashes. Moreover, the timeout is also used by
the controller to calculate a car's reaction time, that is for how long (and for
what distance) it will continue riding at a certain speed until it starts
braking (if necessary) - for humans, there isn't a network delay to account for,
but rather a cognitive reaction time, usually considered of 1 second.

In \ref{req:distr_scalability}, we already discussed the expected growth of the
system's dimensions, for which the network shouldn't feel burdened. 

\subsubsection{Economy, costs}
There shouldn't be any particular constraints to limit economical costs: once
servers and network interfaces have been set up, in a real life scenario, there
don't seem to be any measures which could heavily impact the budget (also refer
to scalability (\ref{req:distr_scalability}) for the expected system's dimensions).
